\chapter{总结与未来展望}
\label{chap:conclusion-future-outlook}
\chaptercontents

我们做到了！现在已经到达终点。在本章最后，我们将简要回顾读者通过本书已经掌握的
学习目标。我们不把本章写成传统意义上的结论，而是提供我们对大规模并行计算未来的展望，
以及对其发展将如何影响科学与技术未来走向的看法。

\section{回顾目标}
\label{sec:goals-revisited}

正如我们在引言中所说，我们的首要目标是教会读者如何为大规模并行处理器编程。我们承诺，
只要建立正确的直觉并采用正确的方法，这件事就会变得容易。具体而言，我们承诺重点培养
计算思维技能，使读者能够以适合并行计算的方式思考问题。我们还承诺关注限制并行应用性能
的因素，并提供系统化的方法，通过优化代码来克服这些因素。

我们分三个步骤兑现了这些承诺。第一步中，第 2 章至第 6 章介绍并行计算和 CUDA C++ 的
基本概念，以及使用 CUDA 编写大规模并行代码时需要考虑的关键性能因素。这些章节还介绍了
理解高性能并行编程中必须应对的硬件限制所需的相关计算机体系结构概念。掌握这些知识后，
开发者就能有把握地编写并行代码，并推理不同线程组织方式、循环结构、数据布局和编码风格
的相对优劣。我们以一份常见优化技术清单结束这部分内容，并在本书后续部分用这些技术优化
各种并行模式和应用的性能。

第二步中，我们介绍了主要的并行模式（第 7 章至第 15 章）。事实证明，这些模式对于在许多
应用中引入并行性很有用。这些章节讲解最有用的并行计算模式背后的概念。每种模式都用具体
的代码示例加以说明，也都用于介绍如何克服并行编程中经常遇到的并行化和性能障碍。我们还
利用这些模式展示第一步中介绍的优化技术如何应用于各种场景。

第三步中，我们介绍了更多高级模式和应用（第 16 章至第 23 章），以巩固前一步获得的知识
和技能。在这一步中，我们继续应用前一步练习过的优化技术，但更加重视探索不同的问题分解
形式，以实现并行化，并分析不同分解方式及其相应数据结构之间的权衡。本步骤的第 22 章专门
回顾计算思维技能，帮助读者把前面章节学到的概念推广为解决新问题所需的高层次思考。掌握
这些认识后，高性能并行编程就成为一个推理过程，而不再是一门黑魔法。

我们希望读者喜欢这本书，并同意自己现在已经为编写大规模并行计算系统做好了充分准备。

\section{未来展望}
\label{sec:future-outlook}

自 2007 年第一款支持 CUDA 的 GPU——G80——问世以来，GPU 作为大规模计算设备的能力已经
取得惊人提升：计算吞吐量提高了 638 倍，内存带宽提高了 93 倍，如图 \ref{fig:24-1} 所示。
这些进步推动了科学和工程领域的高性能计算（HPC）、人工智能和数据分析取得巨大进展，也对
金融、制造和医学等许多垂直领域产生了重要影响。例如，正如第 19 章和第 20 章所述，GPU
从超大规模数据集中开展深度学习，引发了图像识别、语音识别和视频分析领域的革命。

\begin{figure}[htbp]
  \centering
  \renewcommand{\arraystretch}{1.22}
  \small
  \begin{tabularx}{0.98\linewidth}{>{\raggedright\arraybackslash}X
      >{\centering\arraybackslash}X >{\centering\arraybackslash}X
      >{\centering\arraybackslash}X}
    \hline
    & \shortstack{\textbf{G80}\\\textbf{(2007)}} & \shortstack{\textbf{B200}\\\textbf{(2025)}} & \shortstack{\textbf{B200/G80}\\\textbf{比值}} \\
    \hline
    \textbf{计算吞吐量} & \shortstack{0.345 TFLOPS\\(FP32)} & \shortstack{80 TFLOPS /\\2200 TFLOPS\\(FP32 / Tensor Core\\FP32)} & 232 倍 / 638 倍 \\
    \hline
    \textbf{内存带宽} & 86.4 GB/s & 8,000 GB/s & 93 倍 \\
    \hline
    \textbf{内存容量} & 1.5 GB & 192 GB & 128 倍 \\
    \hline
    \textbf{PCIe 带宽} & \shortstack{8 GB/s\\(Gen2 x16，\\每个方向)} & \shortstack{64 GB/s\\(Gen5 x16，\\每个方向)} & 8 倍 \\
    \hline
    \textbf{NVLink 带宽} & 不适用 & 1,800 TB/s & -- \\
    \hline
  \end{tabularx}
  \caption{从 G80 到 B200 的 18 年对比。依据原书图 24.1 重绘；数值和单位按底本保留。}
  \label{fig:24-1}
\end{figure}

\noindent\textbf{译者注：}底本图 24.1 同时列出 G80 的 0.345 TFLOPS、B200 的
2200 TFLOPS，以及二者的比值 638 倍；按前两项计算，$2200/0.345\approx6377$，
与 638 倍不一致，正文中的 638 倍也沿用了该比值。本译稿保留底本原始数值和比值，
并标明这一算术不一致。

虽然单个 GPU 的性能增长极其显著，但自 2017 年以来，通过网络通信协作完成单项任务的 GPU
数量增长得更加惊人。训练大语言模型所需的极高计算复杂度，促使领先科技公司在每项训练任务
中使用数万甚至数十万块 GPU。这些领先科技公司被称为\textbf{超大规模厂商（hyperscaler）}，
因为它们的数据中心规模极其庞大。数据中心中协同工作的 GPU 数量爆炸式增长，推动网络硬件
和软件快速发展。例如，第 23 章介绍了 NCCL；这是一种在超大规模厂商的数据中心训练大语言
模型时得到广泛使用的新型通信库。

自 2010 年本书第 1 版出版以来，并行计算领域也以惊人的速度发展。可用可扩展算法解决的
问题范围显著扩大。GPU 最初主要用于规则的稠密矩阵计算和蒙特卡洛方法，但其应用很快扩展
到稀疏方法和图计算。在许多领域，算法也取得了快速进展。并行模式、高级模式和应用章节中
介绍的一些算法，代表了近年来的重要进展。

有些人自然会怀疑，并行计算的快速发展是否已经走到尽头。从各种迹象来看，答案明确是否定的。
我们很可能正处于并行计算革命的中期。过去三十年计算技术的惊人进步，已经触发了产业界的
范式转变。过去，主要创新通常由物理仪器推动，并由计算设备提供辅助；如今，主要创新则由
计算推动，并由物理仪器提供辅助。

例如，二十年前，GPS 改变了我们的驾驶方式。GPS 主要依靠卫星信号感知，再由计算方法确定
两地之间的最短路径。近来，手机地图应用中的用户上报信息，以及云端数据服务和计算能力，
进一步让驾驶者能够依据交通状况改变路线。如今，汽车行业最令人兴奋的革命是自动驾驶汽车，
它主要依靠机器学习计算方法，并由物理传感器提供辅助。

再以医学为例，MRI 和 PET 在过去二十年中改变了医学。这些技术主要依靠电磁和光学传感器，
并由计算图像重建方法提供辅助；它们让医生无需手术就能看到人体内部的病变。如今，医学正
经历个性化医学的革命，这场革命主要由深度学习模型和计算基因组学方法推动，并由医学和基因
组学仪器提供辅助。

再举一个例子，半导体产业过去依靠物理光源的进步，并由计算方法辅助执行设计规则，以在制造
过程中推动器件特征尺寸缩小。如今，物理光源的进步实际上已经停止。特征尺寸缩小主要由计算
设计的光刻掩模推动：光刻掩模协调光波干涉，从而在芯片上形成极其精确的蚀刻图案。

类似的范式转变也发生在内容创作、药物发现以及天气和气候预报等许多领域。计算已经成为我们
社会几乎所有令人振奋的创新的主要驱动力。这造成了对更快计算系统的永不满足的需求。正如
第 1 章所讨论的，并行计算是提高计算性能的唯一可行途径。这种强劲需求将继续推动产业界
创新，创造更强大的并行计算设备。最有潜力的改进方向之一，是提高访问存储数据时的并行度；
历史上，访问存储数据所采用的并行度一直很低。这一进步可以加速数据库和其他数据密集型应用，
弥补它们在创新方面相对滞后的局面。大幅提高从海量存储数据中提取相关部分的精度和速度，
很可能会激发整整一代我们目前甚至无法想象的应用。

总而言之，我们正处在计算的黄金时代。产业界将继续招募并奖励高技能的并行程序员。读者的
工作将为自己选择的领域带来真正的改变。

\begin{center}
  \large\textit{享受这段旅程！}
\end{center}
